% !TeX program = lualatex
\documentclass[a4paper,twocolumn]{ltjsarticle}
% a4jでは画像挿入時に問題が起こる場合があるため，a4paperを使用する．

% ===== 基本 =====
\usepackage[
    a4paper,
    left=2.0cm,
    right=2.0cm,
    top=2.0cm,
    bottom=2.0cm
]{geometry}
\usepackage{amsmath,amssymb,mathtools,bm}
\usepackage{graphicx}
\usepackage{xcolor}
\usepackage[unicode,hypertexnames=false]{hyperref}
\usepackage{titlesec}
\usepackage{wrapfig}
\usepackage{enumitem}
\usepackage{listings}

% ===== TikZ・図 =====
\usepackage{tikz}
\usetikzlibrary{positioning}
\usepackage{pgfplots}
\pgfplotsset{compat=1.18}

% ===== 化学 =====
\usepackage[version=4]{mhchem}
\usepackage{chemfig}

% ===== 回路 =====
\usepackage{circuitikz}

% ===== アルゴリズム =====
\usepackage{algorithm}
\usepackage{algpseudocode}

% ===== ダミーテキスト =====
% 発展例で図の回り込みを確認するために使用する．実際の文書では不要．
\usepackage[base,english]{babel}
\usepackage{lipsum}

% ===== 数学補助 =====
% \numberwithin{equation}{section}
\renewcommand{\theequation}{\arabic{equation}}

\DeclareMathOperator{\sgn}{sgn}
\DeclareMathOperator{\tr}{tr}
\DeclareMathOperator{\diag}{diag}

\newcommand{\dif}{\mathrm{d}}
\newcommand{\ee}{\mathrm{e}}
\newcommand{\ii}{\mathrm{i}}

\newcommand{\pd}[2]{\frac{\partial #1}{\partial #2}}
\newcommand{\pdd}[2]{\frac{\partial^2 #1}{\partial #2^2}}
\newcommand{\dd}[2]{\frac{\mathrm{d} #1}{\mathrm{d} #2}}

\newcommand{\vect}[1]{\bm{#1}}
\newcommand{\tensor}[1]{\bm{#1}}

% 画像がまだ配置されていない場合も，文書全体はコンパイルできるようにする．
% 指定した画像が存在すれば，通常どおりその画像を表示する．
\newcommand{\practiceimage}[2][\linewidth]{%
  \IfFileExists{#2}{%
    \includegraphics[width=#1]{#2}%
  }{%
    \fbox{%
      \parbox[c][28mm][c]{\dimexpr#1-2\fboxsep-2\fboxrule\relax}{%
        \centering\small
        画像ファイル\\未配置%
      }%
    }%
  }%
}

% ===== 見出し設定 =====
\titleformat{\section}[block]
  {\normalfont\Large\bfseries}
  {\thesection}{1em}{}

\titleformat{\subsection}[block]
  {\normalfont\large\bfseries}
  {\thesubsection}{1em}{}

\titleformat{\subsubsection}[block]
  {\normalfont\normalsize\bfseries}
  {\thesubsubsection}{1em}{}

\titlespacing*{\section}{0pt}{2.0ex plus 0.5ex minus 0.2ex}{0.8ex}
\titlespacing*{\subsection}{0pt}{1.5ex plus 0.4ex minus 0.2ex}{0.5ex}

\lstset{
    language=C,
    basicstyle=\ttfamily\small,
    numbers=left,
    numberstyle=\tiny\color{gray},
    frame=single,
    rulecolor=\color{black},
    backgroundcolor=\color{gray!10},
    breaklines=true,
    showstringspaces=false,
    columns=fullflexible,
    keywordstyle=\color{blue},
    commentstyle=\color{green!50!black},
    stringstyle=\color{red},
    tabsize=4
}

\setlist[itemize]{topsep=4pt,itemsep=2pt,leftmargin=*}
\setlist[enumerate]{topsep=4pt,itemsep=2pt,leftmargin=*}
\setlength{\columnsep}{8mm}
\setlength{\emergencystretch}{1em}
\raggedbottom

\title{%
単振り子の運動とエネルギー
% ここに自分で新しいテーマを付けてみる．テーマは自由である．
}
% ここを自分の名前に変更する．
\author{慎研究室：塩月}
% \date{\today}とすると，コンパイルした日付が入る．
\date{\today}

\begin{document}

\maketitle

% 演習3：Abstractの挿入
% \begin{abstract} ... \end{abstract}を\maketitleの後へ追加する．
\begin{abstract}
本稿では，単振り子の運動を題材として，運動方程式と力学的エネルギーの関係を簡単に整理する．
また，微小振動の近似から振り子の周期を導出し，LaTeXによる数式記述の基本例を示す．
\end{abstract}

\section*{演習1：動作確認}
GitHubからこの「\texttt{LaTeX\_practice.tex}」をダウンロードし，
自分のPCでコンパイルしてみましょう．
正しくコンパイルできれば，PDFが生成されます．

\subsubsection*{解答例}
LuaLaTeXでコンパイルし，PDFが正しく生成されることを確認した．

\section*{演習2：テーマの変更}
タイトルと著者名を自分のものに変更し，\texttt{date}には\verb|\today|を入れて，
コンパイルしてみましょう．

\subsubsection*{解答例}
テーマを「単振り子の運動とエネルギー」，著者名を「山田太郎」とし，
日付には\verb|\today|を指定した．

\section*{演習3：Abstractの挿入}
Abstractを挿入することで，論文の概要を簡潔にまとめることができます．

\begin{verbatim}
\begin{abstract}
ここにAbstractを記述します．
\end{abstract}
\end{verbatim}

実際に，このコードを\verb|\maketitle|の後へ挿入してみましょう．

\subsubsection*{解答例}
文書冒頭に，単振り子の運動方程式と周期について説明するAbstractを挿入した．

\section*{演習4：数式の挿入}
LaTeXでは，数式を美しく記述することができます．
数式の記述には，インライン数式とディスプレイ数式の2種類があります．

\subsection*{インライン数式}
円の面積は $S=\pi r^2$ で表される．

オイラーの公式は $e^{i\pi}+1=0$ である．

関数 $f(x)=x^2+2x+1$ を考える．

\subsection*{ディスプレイ数式}
\begin{equation}
\int_{0}^{1} x^2\,\dif x=\frac{1}{3}
\label{eq:integral}
\end{equation}

\begin{equation}
A=
\begin{pmatrix}
1 & 2\\
3 & 4
\end{pmatrix}
\label{eq:matrix}
\end{equation}

\begin{equation}
\lim_{n\to\infty}
\left(1+\frac{1}{n}\right)^n=e
\label{eq:euler}
\end{equation}

ディスプレイ数式は，\verb|\[ ... \]|で囲むと番号なしで表示できます．
番号を付けたい場合は，\verb|equation|環境を使用します．
実際に，インライン数式とディスプレイ数式をそれぞれ使って，
自分で数式を挿入してみましょう．

% ここに自分で数式を挿入する．
\subsubsection*{解答例}
長さ $\ell$ の単振り子が微小振動するとき，角度 $\theta$ は
$\sin\theta\simeq\theta$ と近似できる．
このとき，運動方程式は
\begin{equation}
\pdd{\theta}{t}+\frac{g}{\ell}\theta=0
\label{eq:pendulum}
\end{equation}
となる．したがって，単振り子の周期 $T$ は
\begin{equation}
T=2\pi\sqrt{\frac{\ell}{g}}
\label{eq:period}
\end{equation}
で与えられる．

\section*{演習5：図の挿入}
LaTeXでは，図を挿入することもできます．

\begin{figure}[htbp]
    \centering
    \practiceimage[3cm]{LaTeX_practice_figure_1.jpg}
    \caption{勉強中のキャラクター画像}
    \label{fig:practice}
\end{figure}

画像の挿入は，画像ファイルの場所や拡張子，パッケージ設定などが原因で
エラーになることがあります．
エラーが出ても慌てず，エラーメッセージを読みながら一つずつ原因を確認しましょう．
では，自分で画像を挿入してみましょう．

% ここに画像を挿入する．
% 分からないときは，画像とともに生成AIへ
% 「画像を貼るLaTeXのコードを作成」などと質問してもよい．

\subsubsection*{解答例}
\begin{figure}[htbp]
    \centering
    \includegraphics[width=0.72\linewidth]{LaTeX_practice_figure_2.jpg}
    \caption{校舎の外観}
    \label{fig:school}
\end{figure}
図\ref{fig:school}のように，プロジェクト内の画像を相対パスで挿入した．

\section*{演習6：表の挿入}
LaTeXでは，表を挿入することもできます．

\begin{table}[htbp]
    \centering
    \begin{tabular}{|c|c|c|}
        \hline
        A & B & C \\
        \hline
        1 & 2 & 3 \\
        4 & 5 & 6 \\
        \hline
    \end{tabular}
    \caption{LaTeXで挿入した表の例}
    \label{tab:practice}
\end{table}

では，自分で5行7列の表を挿入してみましょう．

% ここに表を挿入する．
% 分からないときは，表の画像とともに生成AIへ
% 「表を作るLaTeXのコードを作成」などと質問してもよい．

\subsubsection*{解答例}
\begin{table}[htbp]
    \centering
    \resizebox{\columnwidth}{!}{%
        \begin{tabular}{|c|c|c|c|c|c|c|}
            \hline
            No. & $\ell$ (m) & $g$ (m/s$^2$) & $T$ (s) & $\theta_0$ (deg) & $m$ (kg) & 備考 \\
            \hline
            1 & 0.25 & 9.81 & 1.00 & 5  & 0.10 & 短い \\
            2 & 0.50 & 9.81 & 1.42 & 10 & 0.10 & 基準 \\
            3 & 0.75 & 9.81 & 1.74 & 10 & 0.15 & 長い \\
            4 & 1.00 & 9.81 & 2.01 & 15 & 0.20 & 最長 \\
            \hline
        \end{tabular}%
    }
    \caption{単振り子の条件をまとめた5行7列の表}
    \label{tab:pendulum}
\end{table}

\section*{演習7：箇条書き}
\begin{itemize}
    \item LaTeX
    \item Python
    \item GitHub
\end{itemize}

では，自分で箇条書きを作成してみましょう．

% ここに箇条書きを作成する．

\subsubsection*{解答例}
単振り子の周期を調べる際には，次の点に注意する．
\begin{itemize}
    \item 糸の長さを正確に測定する．
    \item 振れ角を十分に小さくする．
    \item 複数回の振動時間を測り，平均を求める．
\end{itemize}

\section*{演習8：ソースコードの挿入}
LaTeXでは，ソースコードもきれいに掲載できます．

\begin{lstlisting}
#include <stdio.h>

int main(void)
{
    for (int i = 0; i < 10; i++) {
        printf("%d\n", i);
    }

    return 0;
}
\end{lstlisting}

では，自分でソースコードを挿入してみましょう．

% ここにソースコードを挿入する．内容は自由である．

\subsubsection*{解答例}
次のCプログラムは，振り子の長さから周期を計算する．
\begin{lstlisting}
#include <math.h>
#include <stdio.h>

int main(void)
{
    const double length = 0.50;
    const double gravity = 9.81;
    const double pi = 3.141592653589793;
    const double period = 2.0 * pi * sqrt(length / gravity);

    printf("Period: %.2f s\n", period);
    return 0;
}
\end{lstlisting}

\section*{演習9：URLの挿入}
URLを挿入することもできます．
OpenAIのホームページはこちらです．

\url{https://openai.com}

では，自分で任意のURLを挿入してみましょう．

% ここにURLを挿入する．

\subsubsection*{解答例}
振り子に関する解説ページを次に示す．

\url{https://en.wikipedia.org/wiki/Pendulum}

\section*{演習10：参考文献}
卒業論文や学術論文では，他者の研究や論文を引用した場合，
必ず参考文献を明記する必要があります．

LaTeXでは，一般的にBibTeXを用いて参考文献を管理します．
参考文献を追加する手順は，大きく分けて次の3つです．

\subsection*{手順1：\texttt{.bib}ファイルを作成する}
まず，\texttt{references.bib}という名前のファイルを作成します．
このファイルには，引用する論文や書籍の情報を記述します．
例えば，次のように記述します．

\begin{lstlisting}[language={}]
@inproceedings{he2016resnet,
  author    = {Kaiming He and others},
  title     = {Deep Residual Learning for Image Recognition},
  booktitle = {CVPR},
  year      = {2016}
}
\end{lstlisting}

Google ScholarやIEEE Xploreでは，「BibTeX」を選択するだけで
この形式を取得できるため，自分で一から入力する必要はほとんどありません．

\subsection*{手順2：本文中で引用する}
引用したい箇所に，次のように記述します．

\begin{verbatim}
\cite{he2016resnet}
\end{verbatim}

例えば，

\begin{verbatim}
深層学習ではResNetが広く利用されている
\cite{he2016resnet}．
\end{verbatim}

と書くと，「深層学習ではResNetが広く利用されている[1]．」のように，
引用番号が自動で付与されます．

\subsection*{手順3：参考文献一覧を表示する}
文書の末尾へ，次の2行を追加します．

\begin{verbatim}
\bibliographystyle{ieeetr}
\bibliography{references}
\end{verbatim}

ここで，

\begin{itemize}
    \item \texttt{ieeetr}：参考文献の表示形式
    \item \texttt{references}：作成した\texttt{references.bib}の指定
\end{itemize}

を表しています．
実際に手順1から手順3までを行うと，本文では次のように表示されます．

\subsubsection*{解答例}
深層学習では，残差学習を導入したResNetが画像認識に広く利用されている
\cite{he2016resnet}．

卒業論文では，参考文献の管理にBibTeXを用いることが一般的です．
今回は演習を省略しますが，研究を進める際には必ず利用する機能です．

\subsection*{発展：2段組で大きな図を挿入する}
卒業論文では，本文を2段組にし，図やグラフのみページ全体の幅で
表示することがあります．

その場合は，\verb|figure|の代わりに\verb|figure*|環境を使用します．
図がページ上部へ自動配置されることを確認するため，後ろにダミーテキストを入れています．

\begin{lstlisting}[language={[LaTeX]TeX}]
\begin{figure*}[t]
    \centering
    \includegraphics[width=0.8\textwidth]{graph.pdf}
    \caption{実験結果}
    \label{fig:result}
\end{figure*}
\end{lstlisting}

\begin{figure*}[t]
    \centering
    \begin{minipage}{0.47\textwidth}
        \centering
        \practiceimage[\linewidth]{LaTeX_practice_figure_1.jpg}
        \par\smallskip
        \textbf{(a)} 勉強中のキャラクター
    \end{minipage}
    \hfill
    \begin{minipage}{0.47\textwidth}
        \centering
        \practiceimage[\linewidth]{LaTeX_practice_figure_2.jpg}
        \par\smallskip
        \textbf{(b)} 校舎の外観
    \end{minipage}
    \caption{2枚の画像を横並びにした例}
    \label{fig:images}
\end{figure*}

\lipsum[2-7]

% 参考文献一覧は，文書の最後に一度だけ出力する．
\clearpage
\bibliographystyle{ieeetr}
\bibliography{references}

\end{document}